\chapter{Workloads: Pods, Deployments, StatefulSets}
\label{ch:workloads}

\section{Pods}

The pod is the unit Kubernetes schedules: one or more containers sharing
a network namespace (one IP, \code{localhost} between them) and volumes.
Almost every pod here is a single container; the multi-container
patterns (sidecars for proxies or log shippers) matter when you meet a
service mesh. Pods are \emph{mortal by design} --- they are never healed,
only replaced, and their IP dies with them. Everything else in Part~V
--- controllers, Services, PVCs --- exists to build durable behavior out
of disposable pods.

\section{Deployments and the rollout}

You never create bare pods. A \emph{Deployment} declares ``N replicas
of this pod template''; it manages a \emph{ReplicaSet} (the counter),
which manages the pods. The three-layer indirection is what makes
updates elegant: change the template (a new image tag), and the
Deployment creates a \emph{new} ReplicaSet, scaling it up while the old
scales down:

\begin{figure}[h]
\centering
\begin{tikzpicture}[node distance=3.5mm and 9mm]
  \node[box] (d) {Deployment\\\scriptsize course-service};
  \node[box, below left=of d] (rs1) {ReplicaSet (old)\\\scriptsize image :a1b2c3d};
  \node[box, below right=of d] (rs2) {ReplicaSet (new)\\\scriptsize image :e4f5a6b};
  \node[boxfill, below=of rs1] (p1) {pod \scriptsize(terminating)};
  \node[boxfill, below=of rs2] (p2) {pod \scriptsize(starting)};
  \draw[flow] (d) -- (rs1); \draw[flow] (d) -- (rs2);
  \draw[flowdash] (rs1) -- (p1); \draw[flow] (rs2) -- (p2);
\end{tikzpicture}
\caption{A rolling update mid-flight: two ReplicaSets, traffic shifting.}
\end{figure}

The rollout only proceeds as new pods become \emph{Ready}
(Chapter~17's probes --- this is where they earn their keep), so a bad
image stops the rollout instead of taking the service down; and
\code{kubectl rollout undo} points back at the previous ReplicaSet,
which is kept around precisely for this. The pipeline's
\code{kubectl rollout status} (Chapter~23) is CI waiting on this
machinery to conclude.

With one replica --- this cluster --- a rolling update still means
start-new-then-stop-old, a brief two-pod overlap costing a moment of
double memory. \code{maxUnavailable}/\code{maxSurge} tune that dance at
scale.

\section{StatefulSets: when identity matters}

For the databases, interchangeability is exactly wrong: whichever pod
comes up must be \emph{the same Postgres} with \emph{the same data}. A
StatefulSet provides what a Deployment refuses to: stable names
(\code{postgres-course-0}, not a random suffix), per-replica PVCs via
\code{volumeClaimTemplates} (Chapter~19) that survive the pod and
re-attach to its successor, and ordered, one-at-a-time replacement.
From \code{postgres-course.yaml}:

\begin{lstlisting}[language=yaml]
kind: StatefulSet
spec:
  serviceName: postgres-course
  replicas: 1
  volumeClaimTemplates:          # a PVC per replica, not per manifest
    - metadata: { name: data }
      spec:
        accessModes: [ReadWriteOnce]
        resources: { requests: { storage: 5Gi } }
\end{lstlisting}

The decision rule: does the \emph{process} hold state the next process
must inherit? Databases yes; Spring services no --- their state lives in
those databases, which is what made them Deployments, and what makes
them scalable and rollable at all. (The JWT statelessness of
Chapter~5 is this same property one layer up.)

\begin{pitfall}[the Deployment-database data loss]
Symptom: someone runs a database as a Deployment with a plain volume; it
works for months; then a node problem creates a \emph{second} replica
pointing at the same files, and the database corrupts. Cause: a
Deployment's contract is ``interchangeable replicas'' and its rolling
update briefly runs old and new \emph{simultaneously} --- fatal when
both open the same data directory. StatefulSet ordering + per-replica
claims exist to make this impossible. The lesson generalizes: choose
the controller by its \emph{contract}, not by which YAML you saw first.
\end{pitfall}

\section{Other workload kinds}

For completeness, the remaining controllers, none yet used here:
\textbf{DaemonSet} (one pod per node --- log collectors, node metrics),
\textbf{Job} (run to completion --- Chapter~7's Flyway-as-CI-step would
be one), \textbf{CronJob} (Jobs on a schedule --- backups, Chapter~19's
missing piece).

\begin{handson}[watch a rollout happen]
\begin{lstlisting}[language=bash]
kubectl -n ts get rs -l app=course-service          # one ReplicaSet
kubectl -n ts set image deployment/course-service \
    course-service=localhost:5000/course-service:1.0.0-SNAPSHOT
kubectl -n ts rollout status deployment/course-service
kubectl -n ts get rs -l app=course-service          # two -- old kept at 0
kubectl -n ts rollout undo deployment/course-service
\end{lstlisting}
The second \code{get rs} is the point: the old ReplicaSet still exists,
scaled to zero --- the undo target, kept by default for the last ten
revisions.
\end{handson}

\begin{quiz}
\begin{enumerate}
  \item Why does a Deployment manage ReplicaSets instead of pods
  directly? What feature falls out of that indirection?
  \item State the contract difference between Deployment and
  StatefulSet in one sentence each, then classify: Eureka server, a
  Redis cache, a Kafka broker, the api-gateway.
  \item During a rolling update, what condition gates old-pod
  termination, and which chapter's machinery provides it?
  \item Explain the exact mechanism by which running Postgres under a
  Deployment invites corruption where a StatefulSet does not.
\end{enumerate}
\end{quiz}
